% Generated by figures/13_14_24_25_granger_tables.R (score: a2) -- do not edit by hand.
% Requires \usepackage{booktabs} and \usepackage{threeparttable}.
\begin{table}[htbp]
\centering
\begin{threeparttable}
\caption{Granger causality between US and Chinese sentiment by topic, data after 2015}
\label{tab:granger-topics-2015}
\small
\begin{tabular}{l rr r}
\toprule
Topic & US as & China as & Months \\
 & response & response &  \\
\midrule
Domestic periphery \& human rights & 0.400 & 0.698 & 42 \\
Multilateral negotiations & \textbf{0.043} & 0.139 & 80 \\
Trade \& finance & \textbf{0.013} & 0.898 & 79 \\
Taiwan & 0.850 & 0.073 & 22 \\
South China Sea & 0.209 & 0.887 & 36 \\
Global health \& epidemics & 0.229 & 0.536 & 22 \\
Russia & \textbf{0.045} & \textbf{0.027} & 69 \\
\bottomrule
\end{tabular}
\begin{tablenotes}[flushleft]
\footnotesize
\item \textit{Notes:} Wald test ($F$) $p$-values of bivariate Granger causality tests with 3 lags of each series. ``US as response'' tests whether lagged Chinese sentiment improves the prediction of US sentiment beyond lagged US sentiment; ``China as response'' tests the reverse. Bold: $p < 0.05$. Sentiment is the mean hybrid score (approach 2) of the documents in each period. Each series is the monthly mean score of the documents matching a topic's keywords (full US and Chinese data). Months in which either country has no document on the topic are dropped before lagging; ``Months'' counts the months left. Only months after 1 January 2015 are used.
\end{tablenotes}
\end{threeparttable}
\end{table}
